                                    8th IMC 2001
                                   July 19 - July 25
                                Prague, Czech Republic

    Second day

   Problem 1.
   Let r, s ≥ 1 be integers and a0 , a1 , . . . , ar−1 , b0 , b1 , . . . , bs−1 be real non-
negative numbers such that

(a0 + a1 x + a2 x2 + . . . + ar−1 xr−1 + xr )(b0 + b1 x + b2 x2 + . . . + bs−1 xs−1 + xs ) =

                          1 + x + x2 + . . . + xr+s−1 + xr+s .
Prove that each ai and each bj equals either 0 or 1.

   Solution. Multiply the left hand side polynomials. We obtain the following
equalities:
                     a0 b0 = 1, a0 b1 + a1 b0 = 1, . . .
Among them one can find equations

                            a0 + a1 bs−1 + a2 bs−2 + . . . = 1
and

                           b0 + b1 ar−1 + b2 ar−2 + . . . = 1.
    From these equations it follows that a0 , b0 ≤ 1. Taking into account that
a0 b0 = 1 we can see that a0 = b0 = 1.
    Now looking at the following equations we notice that all a’s must be less
than or equal to 1. The same statement holds for the b’s. It follows from
a0 b1 + a1 b0 = 1 that one of the numbers a1 , b1 equals 0 while the other one must
be 1. Follow by induction.

    Problem 2.
            √                 q   p                             2bn
    Let a0 = 2, b0 = 2, an+1 = 2 − 4 − a2n , bn+1 =             p         .
                                                            2 + 4 + b2n
    a) Prove that the sequences (an ), (bn ) are decreasing and converge to 0.
    b) Prove that the sequence (2n an ) is increasing, the sequence (2n bn ) is de-
creasing and that these two sequences converge to the same limit.
    c) Prove that there is a positive constant C such that for all n the following
                                    C
inequality holds: 0 < bn − an < n .
                                   8
                                   p     √       √
p Solution.     Obviously a2 = 2 − 2 < 2. Since the function f (x) =
       √
  2 − 4 − x2 is increasing on the interval [0, 2] the inequality a1 > a2 implies
that a2 > a3 . Simple induction ends the  proof of monotonicity of (an ). In the
                                                                        2x
same way we prove that (bn ) decreases just notice that g(x) =         √         =
                                                                  2 +   4 + x2
          p          
2/ 2/x + 1 + 4/x2 . It is a matter of simple manipulation to prove that
2f (x) > x for all x ∈ (0, 2), this implies that the sequence (2n an ) is strictly

                                             1
increasing. The inequality 2g(x) < x for x ∈ (0, 2) implies that the sequence
                                                                           4b2
(2n bn ) strictly decreases. By an easy induction one can show that a2n = 4+bn2
                                                                              n
for positive integers n. Since the limit of the decreasing sequence (2n bn ) of
positive numbers is finite we have

                                            4 · 4n b2n
                         lim 4n a2n = lim              = lim 4n b2n .
                                             4 + b2n

We know already that the limits lim 2n an and lim 2n bn are equal. The first
of the two is positive because the sequence (2n an ) is strictly increasing. The
existence of a number C follows easily from the equalities

                            4n+1 b2n  n               (2n bn )4 1       1
2n bn − 2n an = 4n b2n −          2
                                      / 2 bn + 2 n a n =         ·  ·
                            4 + bn                       4 + b2n 4n 2n (bn + an )

and from the existence of positive limits lim 2n bn and lim 2n an .
   Remark.       The last problem may be solved in a much simpler way by
someone who is able to make use of sine and cosine. It is enough to notice that
            π                     π
an = 2 sin n+1 and bn = 2 tan n+1 .
          2                     2
   Problem 3.
   Find the maximum number of points on a sphere of radius 1 in Rn√such that
the distance between any two of these points is strictly greater than 2.

   Solution. The unit sphere in Rn is defined by
                                                n
                        (                               )
                                                X
                                             n     2
                Sn−1 = (x1 , . . . , xn ) ∈ R |   xk = 1 .
                                                           k=1

   The distance between the points X = (x1 , . . . , xn ) and Y = (y1 , . . . , yn ) is:
                                              n
                                              X
                               d2 (X, Y ) =         (xk − yk )2 .
                                              k=1

We have

                           √
              d(X, Y ) >       2 ⇔ d2 (X, Y ) > 2
                                   Xn        Xn          n
                                                         X
                                 ⇔     x2k +     yk2 + 2   xk y k > 2
                                       k=1           k=1            k=1
                                       Xn
                                  ⇔           xk y k < 0
                                       k=1

   Taking account of the symmetry of the sphere, we can suppose that

                                  A1 = (−1, 0, . . . , 0).
                   n
                   P
   For X = A1 ,          xk yk < 0 implies y1 > 0, ∀ Y ∈ Mn .
                   k=1
   Let X = (x1 , X), Y = (y1 , Y ) ∈ Mn \{A1 }, X, Y ∈ Rn−1 .


                                               2
   We have
                n
                X                               n−1
                                                X                      n−1
                                                                       X
                      xk y k < 0 ⇒ x 1 y 1 +            xk y k < 0 ⇔          x0k yk0 < 0,
                k=1                             k=1                     k=1
   where
                                      xk           y
                               x0k = pP 2 , yk0 = pPk 2 .
                                         xk          yk
   therefore

                           (x01 , . . . , x0n−1 ), (y10 , . . . , yn−1
                                                                   0
                                                                       ) ∈ Sn−2
                   n
                   P
   and verifies          xk yk < 0.
                  k=1
   If an is the search number of points in Rn we obtain an ≤ 1 + an−1 and
a1 = 2 implies that an ≤ n + 1.
   We show that an = n + 1, giving an example of a set Mn with (n + 1)
elements satisfying the conditions of the problem.
                        A1 = (−1, 0, 0, 0, . . . , 0, 0) 
                        A2 = n1 , −c1 , 0, 0, . . . , 0, 0            
                                  1
                        A3 =        , 1 · c1 , −c2 , 0, . . . , 0, 0
                                 n n−1                                           
                                  1    1          1
                        A4 =      n , n−1 · c1 , n−1 · c2 , −c3 , . . . , 0, 0
                                                                               
                         1
           An−1 = n1 , n−1           1
                             · c1 , n−2          1
                                         · c2 , n−3  · c3 , . . . , −cn−2 , 0
                                                                                   
           An = n1 , n−1
                      1           1
                          · c1 , n−2          1
                                      · c1 , n−3  · c3 , . . . , 12 · cn−2 , −cn−1
                                                                                    
           An+1 = n1 , n−1
                         1           1
                             · c1 , n−2          1
                                         · c2 , n−3  · c3 , . . . , 12 · cn−2 , cn−1
   where                 s                           
                                    1               1
                 ck =         1+             1−           , k = 1, n − 1.
                                    n             n−k+1
           n                                  n
                xk yk = − n1 < 0 and                x2k = 1, ∀X, Y ∈ {A1 , . . . , An+1 } .
           P                                  P
We have
          k=1                                k−=1
    These points are on the unit sphere in Rn and the distance between any two
points is equal to
                                    √             √
                                      r
                                             1
                               d = 2 1 + > 2.
                                            n
    Remark. For n = 2 the points form an equilateral triangle in the unit
circle; for n = 3 the four points from a regular tetrahedron and in Rn the points
from an n dimensional regular simplex.

    Problem 4.
    Let A = (ak,` )k,`=1,...,n be an n × n complex matrix such that for each
m ∈ {1, . . . , n} and 1 ≤ j1 < . . . < jm ≤ n the determinant of the matrix
(ajk ,j` )k,`=1,...,m is zero. Prove that An = 0 and that there exists a permutation
σ ∈ Sn such that the matrix
                                        (aσ(k),σ(`) )k,`=1,...,n

                                                    3
has all of its nonzero elements above the diagonal.

    Solution. We will only prove (2), since it implies (1). Consider a directed
graph G with n vertices V1 , . . . , Vn and a directed edge from Vk to V` when
ak,` 6= 0. We shall prove that it is acyclic.
    Assume that there exists a cycle and take one of minimum length m. Let
j1 < . . . < jm be the vertices the cycle goes through and let σ0 ∈ Sn be a
permutation such that ajk ,jσ0 (k) 6= 0 for k = 1, . . . , m. Observe that for any
other σ ∈ Sn we have ajk ,jσ(k) = 0 for some k ∈ {1, . . . , m}, otherwise we would
obtain a different cycle through the same set of vertices and, consequently, a
shorter cycle. Finally

                               0 = det(ajk ,j` )k,`=1,...,m


                         m                                           m
         = (−1)sign σ0                                  (−1)sign σ
                         Y                      X                    Y
                               ajk ,jσ0 (k) +                              ajk ,jσ(k) 6= 0,
                         k=1                    σ6=σ0                k=1

which is a contradiction.
    Since G is acyclic there exists a topological ordering i.e. a permutation
σ ∈ Sn such that k < ` whenever there is an edge from Vσ(k) to Vσ(`) . It is easy
to see that this permutation solves the problem.

   Problem 5. Let R be the set of real numbers. Prove that there is no
function f : R → R with f (0) > 0, and such that

                  f (x + y) ≥ f (x) + yf (f (x)) for all x, y ∈ R.

    Solution. Suppose that there exists a function satisfying the inequality. If
f (f (x)) ≤ 0 for all x, then f is a decreasing function in view of the inequalities
f (x + y) ≥ f (x) + yf (f (x)) ≥ f (x) for any y ≤ 0. Since f (0) > 0 ≥ f (f (x)),
it implies f (x) > 0 for all x, which is a contradiction. Hence there is a z such
that f (f (z)) > 0. Then the inequality f (z + x) ≥ f (z) + xf (f (z)) shows that
 lim f (x) = +∞ and therefore lim f (f (x)) = +∞. In particular, there exist
x→∞                                  x→∞
x, y > 0 such that f (x) ≥ 0, f (f (x)) > 1, y ≥ f (fx+1
                                                     (x))−1 and f (f (x + y + 1)) ≥ 0.
Then f (x + y) ≥ f (x) + yf (f (x)) ≥ x + y + 1 and hence
                                                              
 f (f (x + y)) ≥ f (x + y + 1) + f (x + y) − (x + y + 1) f (f (x + y + 1)) ≥
                ≥ f (x + y + 1) ≥ f (x + y) + f (f (x + y)) ≥
                ≥ f (x) + yf (f (x)) + f (f (x + y)) > f (f (x + y)).

This contradiction completes the solution of the problem.




                                                4
   Problem 6.
   For each positive integer n, let fn (ϑ) = sin ϑ · sin(2ϑ) · sin(4ϑ) · · · sin(2n ϑ).
   For all real ϑ and all n, prove that
                                          2
                               |fn (ϑ)| ≤ √ |fn (π/3)|.
                                           3

   Solution.
       √       We prove that g(ϑ) = | sin ϑ|| sin(2ϑ)|1/2 attains its maximum
value ( 3/2)3/2 at points 2k π/3 (where k is a positive integer). This can be
seen by using derivatives or a classical bound like

                                   √ q                                                2
                                     2 4                                    √
  |g(ϑ)| = | sin ϑ|| sin(2ϑ)|1/2 = √
                                   4
                                         | sin ϑ| · | sin ϑ| · | sin ϑ| · |   3 cos ϑ|
                                     3
                        √                                       √ !3/2
                          2 3 sin2 ϑ + 3 cos2 ϑ                   3
                      ≤ √   ·                   =                      .
                        4
                          3          4                           2
   Hence

                                                                                        1−E/2
  fn (ϑ)        g(ϑ) · g(2ϑ)1/2 · g(4ϑ)3/4 · · · g(2n−1 ϑ)E      sin(2n ϑ)
          =                                                    ·
 fn (π/3)   g(π/3) · g(2π/3)1/2 · g(4π/3)3/4 · · · g(2n−1 π/3)E sin(2n π/3)
                                    1−E/2                     1−E/2
                       sin(2n ϑ)
                                                
                                                        1               2
                  ≤                         ≤       √                  ≤√ .
                      sin(2n π/3)                       3/2              3
   where E = 32 (1 − (−1/2)n ). This is exactly the bound we had to prove.




                                            5
